\section{Introduction}

Theta functions are special functions of several complex variables that are central to the theory of modular forms, elliptic functions, and abelian varieties. Gauss's work on the arithmetic-geometric mean (AGM) and theta functions laid the foundation for much of modern number theory and mathematical physics.

In this chapter, we cover:
\begin{itemize}
    \item Jacobi theta functions $\vartheta_1, \vartheta_2, \vartheta_3, \vartheta_4$
    \item Modular transformations and the modular group $SL(2,\Z)$
    \item Theta functions of quadratic forms
    \item The Dedekind eta function
    \item Applications to number theory (sums of squares, partition functions)
\end{itemize}

\section{Jacobi Theta Functions}

\begin{definition}[Jacobi Theta Functions]
For $z \in \C$, $\tau \in \mathbb{H}$ (upper half-plane), $q = e^{\pi i \tau}$:
\begin{align*}
\vartheta_1(z|\tau) &= -i \sum_{n=-\infty}^{\infty} (-1)^n q^{(n+1/2)^2} e^{2\pi i (n+1/2)z} \\
\vartheta_2(z|\tau) &= \sum_{n=-\infty}^{\infty} q^{(n+1/2)^2} e^{2\pi i (n+1/2)z} \\
\vartheta_3(z|\tau) &= \sum_{n=-\infty}^{\infty} q^{n^2} e^{2\pi i n z} \\
\vartheta_4(z|\tau) &= \sum_{n=-\infty}^{\infty} (-1)^n q^{n^2} e^{2\pi i n z}
\end{align*}
\end{definition}

\begin{theorem}[Jacobi's Triple Product Identity]
\[
\vartheta_1(z|\tau) = 2 q^{1/4} \sin(\pi z) \prod_{n=1}^{\infty} (1 - q^{2n})(1 - 2 q^{2n} \cos(2\pi z) + q^{4n}).
\]
\end{theorem}

\section{Modular Transformations}

The modular group $\Gamma = SL(2,\Z)$ acts on $\mathbb{H}$ by $\gamma \tau = \frac{a\tau + b}{c\tau + d}$ for $\gamma = \begin{pmatrix} a & b \\ c & d \end{pmatrix}$.

\begin{theorem}[Modular Transformation of $\vartheta_3$]
\[
\vartheta_3(0|-1/\tau) = \sqrt{-i\tau}\, \vartheta_3(0|\tau).
\]
More generally, theta functions transform with specific factors (the theta multiplier system).
\end{theorem}

\section{Theta Functions of Quadratic Forms}

For a positive definite quadratic form $Q(x) = x^T A x$ with $A \in M_n(\Z)$,
\[
\Theta_Q(z) = \sum_{x \in \Z^n} e^{\pi i Q(x) z}.
\]
This is a modular form of weight $n/2$.

\begin{example}[Sum of Squares]
$r_k(n) =$ number of ways to write $n$ as sum of $k$ squares.
The generating function is $\Theta_{I_k}(z)^k$ where $I_k$ is the $k\times k$ identity.
\end{example}

\section{Python Implementation}

In \texttt{python/theta.py}:

\begin{lstlisting}[style=python, caption={Key functions in python/theta.py}]
theta1(z, tau), theta2(z, tau), theta3(z, tau), theta4(z, tau)
theta_null(tau)                   # theta(0|tau)
modular_transform(theta, tau)     # Transform under SL(2,Z)
theta_quadratic_form(A, z)        # Theta of quadratic form
sum_of_squares(k, N)              # r_k(n) for n <= N
dedekind_eta(tau)                 # eta function
\end{lstlisting}

\section{Numerical Examples}

\begin{example}[Theta Null Values]
\begin{lstlisting}[style=python]
>>> from theta import theta_null
>>> import cmath
>>> for tau in [1j, 2j, 0.5+1j]:
...     t = theta_null(tau)
...     print(f"tau={tau}: theta3(0)={t:.6f}")
tau=1j: theta3(0)=1.772637
tau=2j: theta3(0)=1.180341
tau=0.5+1j: theta3(0)=1.355623-0.301032j
\end{lstlisting}
\end{example}

\begin{example}[Sums of Squares]
\begin{lstlisting}[style=python]
>>> from theta import sum_of_squares
>>> r4 = sum_of_squares(4, 20)
>>> for n in range(1, 11):
...     print(f"r_4({n}) = {r4[n]}")
r_4(1) = 8
r_4(2) = 24
r_4(3) = 32
r_4(4) = 24
r_4(5) = 48
...
\end{lstlisting}
\end{example}

\section{Exercises}

\begin{exercise}
Implement the modular transformation of $\vartheta_3$ and verify numerically for various $\tau$.
\end{exercise}

\begin{exercise}
Use theta functions to prove Jacobi's four-square theorem: $r_4(n) = 8 \sum_{d|n, 4\nmid d} d$.
\end{exercise}

\begin{exercise}
Implement the Dedekind eta function and verify its modular transformation $\eta(-1/\tau) = \sqrt{-i\tau}\, \eta(\tau)$.
\end{exercise}

\begin{exercise}
Compute the Fourier expansion of the $j$-invariant using $\eta$ and $\vartheta$ functions.
\end{exercise}

\section{References}

\begin{itemize}
    \item \cite{gauss-works} -- Gauss's works on theta and AGM
    \item \cite{whittaker-watson} -- Whittaker \& Watson, Ch. 21 (Theta functions)
    \item \cite{serre-modular} -- Serre, \textit{A Course in Arithmetic}, Ch. 3
    \item \cite{mumford-tata} -- Mumford, \textit{Tata Lectures on Theta}
\end{itemize}